\documentclass[french]{book}
\usepackage{teach}
\usepackage{verbatim}
\usepackage{amsmath,amsfonts,amssymb}
\usepackage{pgfplots}
\usepackage{mwe}
\RequirePackage{graphicx}
\RequirePackage{ifpdf}
 \ifpdf
   \DeclareGraphicsRule{*}{mps}{*}{}
 \fi


  \usepackage{fancyhdr}

  \usepackage{amsmath}
  \usepackage{amssymb}
  \usepackage{amsfonts}
  \usepackage{amsthm}
  \usepackage{mathrsfs}

  \usepackage{array}
  \usepackage{multicol}
  \setlength{\multicolsep}{3pt}

  \usepackage{graphicx}
  \graphicspath{{Figures/}}
  \usepackage{pstricks,pst-plot,pst-text,pst-tree,pst-eucl}
  \usepackage[all]{xy}

  \usepackage{fancybox}
  \usepackage{color}

%  \usepackage[frenchb]{babel}
  \usepackage{hyperref}

\usepackage{subfig}
\pgfplotsset{compat=newest}
\usepackage[all]{xy}
%\usepackage{geometry}
%\geometry{top=1cm, bottom=1cm, left=2cm, right=2cm}
\usepackage[latin1]{inputenc}
\usepackage[french]{babel}
\redefinecolor{thm}{title}{bgcolor}{alizarin}
\redefinecolor{prop}{title}{bgcolor}{meatbrown}
\redefinecolor{defn}{title}{bgcolor}{olivine}
\definecolor{alizarin}{rgb}{0.82, 0.1, 0.26}
\definecolor{meatbrown}{rgb}{0.9, 0.72, 0.23}
\definecolor{olivine}{rgb}{0.6, 0.73, 0.45}
\usepackage{mathrsfs}
%%
\newcommand{\R}{\mathbb {R}}
%\newcommand{\C}{\mathds {C}}
\newcommand{\N}{\mathbb {N}}
\newcommand{\Z}{\mathbb {Z}}
\newcommand{\Q}{\mathbb {Q}}
\newcommand{\D}{\mathbb {D}}
%%
\newcommand{\se}{\geq}
\newcommand{\ie}{\leq}
\newcommand{\tm}{\times}
%\newcommand{\ell}{l}
%\newcommand{\ell'}{l'}
%%
\newcommand{\ds}{\displaystyle}
\newcommand{\limn}{\lim_{n\rightarrow +\infty}}
\newcommand{\limo}[1][x]{\ds\lim_{#1\rightarrow 0}}
\newcommand{\limite}[2][x]{\ds\lim_{#1\rightarrow #2}}
\newcommand{\limpinf}[1][x]{\ds\lim_{#1\rightarrow +\infty}}
\newcommand{\limminf}[1][x]{\ds\lim_{#1\rightarrow -\infty}}
\newcommand{\limiteg}[2][x]{\ds\lim_{#1\xrightarrow{<}#2}}
\newcommand{\limited}[2][x]{\ds\lim_{#1\xrightarrow{>}#2}}
%%
\newcommand{\vect}[1]{\overrightarrow{#1}}
\newcommand{\oij}{(\textit{O},{\overrightarrow{i}},{\overrightarrow{j}})}
%
\newcommand{\cp}[2]{%
        \begin{pmatrix}
        #1\\
        #2
        \end{pmatrix}%
        }
%
\newcommand{\calb}{\mathscr{B}}
\newcommand{\calc}{\mathscr{C}}
\newcommand{\cald}{\mathscr{D}}
\newcommand{\cale}{\mathscr{E}}
\newcommand{\calf}{\mathscr{F}}
\newcommand{\calp}{\mathscr{P}}
\newcommand{\cals}{\mathscr{S}}
\newcommand{\calt}{\mathscr{T}}
%
\newcommand{\suite}[1][u]{\left( #1_{n} \right)_{n\in\N}}
\newcommand{\suitep}[1][u]{\left( #1_{n} \right)_{n\in\N^{*}}}
\usetikzlibrary{arrows}

\newcommand{\poly}[1][a]{#1_0+#1_1x+\cdots+#1_nx^n}


\begin{document}
\tableofcontents

\chapter{Limites}
\section{Limites d'une fonction en l'infini}
\subsection{Limite r\'eelle en l'infini, asymptote horizontale}
\begin{prop}
Soit $f$ une fonction d\'efinie sur un intervalle de la forme $[A~;~+\infty[$.
On dit que $f$ tend vers $\ell$ lorsque $x$ tend vers $+\infty$ si tout intervalle ouvert contenant $\ell$ contient aussi toutes les valeurs de $f(x)$ pour $x$ suffisamment grand.

On note $ \limpinf f(x)=\ell$.
\end{prop}

\begin{minipage}{9cm}
\begin{rem}
On peut d\'efinir de m\^eme $ \limminf f(x)$.
\end{rem}

\begin{exemple}
Soit $f$ la fonction d\'efinie sur $[1~;~+\infty[$ par $f(x)=1+\dfrac{1}{x}$. D\'emontrer que $\limpinf f(x)=1$.
\end{exemple}


Soit $a\in\R$.
\begin{align*}
f(x)\in ]1-a~;~1+a[ &\Leftrightarrow 1-a < 1+\dfrac{1}{x} < 1+a\\
& \Leftrightarrow -a < \dfrac{1}{x} < a\\
& \Leftrightarrow x>\dfrac{1}{a} \quad \text{car $x$ est positif.}
\end{align*}

Ainsi pour $x>\dfrac{1}{a}$, $f(x)\in~ ]1-a~;~1+a[$ donc $\limpinf f(x)=1$.

\end{minipage}
\hfill
\begin{minipage}{5.9cm}
\begin{center}
\includegraphics{CoursLimitesFigures.1}
\end{center}
\end{minipage}

\begin{prop}
Soit $f$ une fonction d\'efinie sur un intervalle de la forme $[A~;~+\infty[$ et $\calc$ sa courbe repr\'esentative dans un rep\`ere.
Si $ \limpinf f(x)=\ell$ ou si  $ \limminf f(x)=\ell$, on dit que la droite d'\'equation $y=\ell$ est asymptote � $\calc$.
\end{prop}

\subsection{Limites usuelles du type $ \frac{1}{x^\alpha}$ en l'infini}
\begin{prop}
$\limpinf \dfrac{1}{x}=0$ ;\hfill
$\limminf \dfrac{1}{x}=0$ ; \hfill $\limpinf \dfrac{1}{x^{2}}=0$ ; \hfill $\limminf \dfrac{1}{x^{2}}=0$ ; \hfill $\limpinf \dfrac{1}{\sqrt{x}}=0$.
\end{prop}

\subsection{Limite infinie en l'infini}
\begin{prop}
Soit $f$ une fonction d\'efinie sur un intervalle de la forme $[A~;~+\infty[$.
On dit que $f$ tend vers $+\infty$ lorsque $x$ tend vers $+\infty$ si tout intervalle de la forme $]M~;~+\infty[$ contient toutes les valeurs de $f(x)$ pour $x$ suffisamment grand.

On note $ \limpinf f(x)=+\infty$.
\end{prop}

\begin{minipage}{9cm}
\begin{rem}
On peut d\'efinir de m\^eme $\limpinf f(x)=-\infty$...
\end{rem}

\begin{exemple}
Soit $f$ la fonction d\'efinie sur $[1~;~+\infty[$ par $f(x)=x^{2}+3$. D\'emontrer que $\limpinf f(x)=+\infty$.
\end{exemple}


Soit $M\in\R$.
\begin{align*}
f(x)>M &\Leftrightarrow x^{2}+3>M \Leftrightarrow x^{2}>M-3\\
& \Leftrightarrow x>\sqrt{M-3} \quad \text{pour $M>3$.}
\end{align*}

Ainsi pour $x>\sqrt{M-3}$, $f(x)>M$ donc $\limpinf f(x)=+\infty$.

\end{minipage}
\hfill
\begin{minipage}{5.9cm}
\begin{center}
\includegraphics{CoursLimitesFigures.2}
\end{center}
\end{minipage}
\subsection{Limites usuelles du type $ x^\alpha $ en l'infini}
\begin{prop}
$\limpinf x=+\infty$ ; \hfill
$\limminf x=-\infty$ ;\hfill $\limpinf x^{2}=+\infty$ ; \hfill $\limminf x^{2}=+\infty$; \hfill $\limpinf \sqrt{x}=+\infty$.
\end{prop}
\newpage

\section{Limite d'une fonction en un point}
\subsection{Limite r\'eelle en un point}
\begin{prop}
Soit $f$ une fonction d\'efinie sur un intervalle $I$ et $a\in I$.
On dit que $f$ tend vers $L$ lorsque $x$ tend vers $a$ si tout intervalle ouvert contenant $L$ contient aussi toutes les valeurs de $f(x)$ pour $x$ suffisamment proche de $a$.

On note $\limite[x]{a}f(x)=L$.
\end{prop}

\begin{rem}
$\limite[x]{a}f(x)=L\Leftrightarrow\limite{h}{0}f(a+h)=L$
\end{rem}

\begin{exemple}
$\limo[x]x=0$ ; \hspace{2em} $\limo[x]x^{2}=0$ ; \hspace{2em} $\limo[x]\dfrac{(1+x)^{2}-1}{x}=2$.
\end{exemple}

\subsection{Limite infinie en un point, asymptote verticale}
\begin{prop}
Soit $f$ une fonction d\'efinie sur un intervalle $I$ et $a\in I$.
On dit que $f$ tend vers $+\infty$ lorsque $x$ tend vers $a$ si tout intervalle de la forme $]M~;~+\infty[$ contient toutes les valeurs de $f(x)$ pour $x$ suffisamment proche de $a$.

On note $\limite[x]{a}f(x)=+\infty$.
\end{prop}

\begin{minipage}{9cm}

\begin{rem}
On peut d\'efinir de m\^eme $\limite[x]{a}f(x)=-\infty$.
\end{rem}

\begin{exemple}
Soit $f$ la fonction d\'efinie sur $\R^{*}$ par $f(x)=\dfrac{1}{x^{2}}$. D\'emontrer que $\limo[x]f(x)=+\infty$.
\end{exemple}


Soit $M\in\R$
\begin{align*}
f(x)>M&\Leftrightarrow\dfrac{1}{x^{2}}>M\Leftrightarrow x^{2}<\dfrac{1}{M}\\
&\Leftrightarrow -\dfrac{1}{\sqrt{M}}<x<\dfrac{1}{\sqrt{M}}
\end{align*}
Ainsi pour $-\dfrac{1}{\sqrt{M}}<x<\dfrac{1}{\sqrt{M}}$, $f(x)>M$ donc $\limo[x]f(x)=+\infty$.

\end{minipage}
\hfill
\begin{minipage}{5cm}
\begin{center}
\includegraphics{CoursLimitesFigures.3}
\end{center}
\end{minipage}
\newpage
\begin{prop}
Soit $f$ une fonction d\'efinie sur un intervalle $I$ et $a\in I$. Soit $\calc$ sa courbe repr\'esentative dans un rep\`ere.

Si $\limite[x]{a}f(x)=+\infty$ ou $\limite[x]{a}f(x)=-\infty$, on dit que la droite d'\'equation $x=a$ est asymptote � $\calc$.
\end{prop}

\subsection{Limites usuelles  du type $ \frac{1}{x^\alpha}$ en 0}
\begin{prop}
$\ds \lim_{_{x>0}^{x \to 0}} \dfrac{1}{x}=+\infty$ ;\hfill $\ds \lim_{_{x<0}^{x \to 0}} \dfrac{1}{x}=-\infty$ ;\hfill
$\limo[x]\dfrac{1}{x^{2}}=+\infty$ ; \hfill $\limo[x]\dfrac{1}{\sqrt{x}}=+\infty$
\end{prop}


\section{Asymptotes obliques}
\begin{prop}
Soit $f$ une fonction d\'efinie sur un intervalle de la forme $[A~;~+\infty[$ et $\calc$ sa courbe repr\'esentative dans un rep\`ere. Soient $a$ et $b$ deux r\'eels avec $a\neq0$.\\

Si $\limpinf (f(x)-(ax+b))=0$ (resp. $\limminf (f(x)-(ax+b))=0$), on dit que la droite d'\'equation $y=ax+b$ est asymptote oblique � la courbe $\calc$ en $+\infty$ (resp. en $-\infty$).
\end{prop}

\begin{minipage}{8cm}

\underline{Interpr\'etation graphique :}\\

Le point $M$ a pour coordonn\'ees $(x;f(x))$ et le point $N$ a pour coordonn\'ees $(x;ax+b)$.

La distance $MN$ est donc \'egale � \[\left|f(x)-(ax+b)\right|\]
Ainsi, si $\limpinf (f(x)-(ax+b))=0$ alors la longueur $MN$ tend vers 0.

La courbe \og se rapproche \fg{} de la droite et tend � \og suivre la direction \fg{} de la droite.
\end{minipage}
\hfill
\begin{minipage}{6.9cm}
\begin{center}
\includegraphics{CoursLimitesFigures.4}
\end{center}
\end{minipage}

\begin{exemple}
Soit $f$ la fonction d\'efinie sur $\R^{*}$ par $f(x)=2x-1+\dfrac{1}{x^{2}}$ et soit $\calc$ sa courbe repr\'esentative. Soit $d$ la droite d'\'equation $y=2x-1$.

D\'emontrer que $d$ est asymptote � $\calc$ en $+\infty$.
\end{exemple}


Pour tout $x \neq 0$, $f(x)-(2x-1)=\dfrac{1}{x^{2}}$.

Or $\limpinf \dfrac{1}{x^{2}}=0$.

Ainsi, $d$ est asymptote � $\calc$ en $+\infty$.

\newpage
\section{Op\'erations sur les limites}
Dans les tableaux suivants, $a$ d\'esigne un nombre r\'eel, $+\infty$ ou $-\infty$.
\subsection{Somme de fonctions}
\begin{center}
\renewcommand{\arraystretch}{2}
\begin{tabular}{|lc|c|c|c|c|c|c|} \hline
Si & $\limite[x]{a}f(x)=$ & $\ell$ & $\ell$ & $\ell$ & $+\infty$ & $+\infty$ & $-\infty$ \\ \hline
et & $\limite[x]{a}g(x)=$ & $\ell'$ & $+\infty$ & $-\infty$ & $+\infty$ & $-\infty$ & $-\infty$ \\ \hline
alors & $\limite[x]{a}f(x)+g(x)=$ & $\ell+\ell'$& $+\infty$ & $-\infty$ & $+\infty$ & ??? & $-\infty$ \\ \hline
\end{tabular}
\end{center}

\begin{rem}
??? signifie que l'on ne peut pas conclure. Cela ne signifie pas que la limite n'existe pas mais simplement qu'on ne peut pas la trouver directement. On parle de \og forme ind\'etermin\'ee \fg.
\end{rem}

\begin{exemple}
D\'eterminer $\limpinf x^{2}+\dfrac{1}{x}$.
\end{exemple}


$\left\lbrace
\begin{array}{l}
\limpinf x^{2}=+\infty \\
\limpinf \dfrac{1}{x}=0
\end{array}
\right.$
donc $\limpinf x^{2}+\dfrac{1}{x}=+\infty$


\subsection{Produit par une constante}
$k$ d\'esigne un nombre r\'eel diff\'erent de 0.
\begin{center}
\renewcommand{\arraystretch}{2}
\begin{tabular}{|lc|c|c|c|c|c|} \hline
Si & $\limite[x]{a}f(x)=$ & $\ell$ & $+\infty$ & $+\infty$ & $-\infty$ & $-\infty$ \\ \hline
et &  &  & $k>0$ & $k<0$ & $k>0$ & $k<0$ \\ \hline
alors & $\limite[x]{a}kf(x)=$ & $k\ell$ & $+\infty$ & $-\infty$ & $-\infty$ & $+\infty$ \\ \hline
\end{tabular}
\end{center}

\subsection{Produit de fonctions}
\begin{center}
\renewcommand{\arraystretch}{2}
\begin{tabular}{|lc|c|c|c|c|c|c|} \hline
Si & $\limite[x]{a}f(x)=$ & $\ell$ & $\ell\neq0$ & 0 & $+\infty$ & $+\infty$ & $-\infty$ \\ \hline
et & $\limite[x]{a}g(x)=$ & $\ell'$ & $\pm\infty$ & $\pm\infty$ & $+\infty$ & $-\infty$ & $-\infty$ \\ \hline
alors & $\limite[x]{a}f(x)\times g(x)=$ & $\ell\times \ell'$ & $\pm\infty$ & ??? & $+\infty$ & $-\infty$ & $+\infty$ \\ \hline
\end{tabular}
\end{center}

\begin{rem}
$\pm\infty$ signifie $+\infty$ ou $-\infty$. Dans le cas de la troisi\`eme ligne, c'est la r\`egle des signes d'un produit qui permet de conclure.
\end{rem}

\begin{exemple}
D\'eterminer $\limpinf \left(\dfrac{1}{x}-1\right)(x^{2}+2)$.
\end{exemple}


$\left\lbrace
\begin{array}{l}
\limpinf \dfrac{1}{x}-1=-1 \\
\limpinf x^{2}+2=+\infty
\end{array}
\right.$
donc $\limpinf \left(\dfrac{1}{x}-1\right)(x^{2}+2)=-\infty$


\subsection{Quotient de fonctions}
\begin{center}
\renewcommand{\arraystretch}{2}
\begin{tabular}{|lc|c|c|c|c|c|c|} \hline
Si & $\limite[x]{a}f(x)=$ & $\ell$ & $\ell$ & $\ell\neq0$ & 0 & $\pm\infty$ & $\pm\infty$ \\ \hline
et & $\limite[x]{a}g(x)=$ & $\ell'\neq0$ & $\pm\infty$ & 0 & 0 & $\ell'$ & $\pm\infty$ \\ \hline
alors & $\limite[x]{a}\dfrac{f(x)}{g(x)}=$ & $\dfrac{\ell}{\ell'}$ & 0 & $\pm\infty$ & ??? & $\pm\infty$ & ??? \\ \hline
\end{tabular}
\end{center}

\begin{exemple}
D\'eterminer $\limite[x]{1}\dfrac{-3x+2}{(x-1)^{2}}$.
\end{exemple}


$\left\lbrace
\begin{array}{l}
\limite[x]{1}-3x+2=-1 \\
\limite[x]{1}(x-1)^{2}=0^{+}
\end{array}
\right.$
donc $\limite[x]{1}\dfrac{-3x+2}{(x-1)^{2}}=-\infty$

\subsection{Exemple d'\'etude d'une forme ind\'etermin\'ee}
Soit $f$ la fonction d\'efinie sur $\R$ par $f(x)=x^{2}-3x+5$.

\'etudier les limites de $f$ en $+\infty$ et $-\infty$.


$\left\lbrace
\begin{array}{l}
\limminf x^{2}=+\infty \\
\limminf -3x+5=+\infty
\end{array}
\right.$
donc $\limminf f(x)=+\infty$

Pour tout $x\neq0$, $f(x)=x^{2}\left(1-\dfrac{3}{x}+\dfrac{5}{x^{2}}\right)$

$\left\lbrace
\begin{array}{l}
\limpinf \dfrac{3}{x}=0\\
\limpinf \dfrac{5}{x^{2}}=0
\end{array}
\right.$
donc $\limpinf 1-\dfrac{3}{x}+\dfrac{5}{x^{2}}=1$

De plus, $\limpinf x^{2}=+\infty$

Ainsi $\limpinf f(x)=+\infty$.

\newpage
\section{Fonction exponentielle : Croissances compar�es}
On sait que l'exponentielle est une fonction qui tend vers $+ \infty$ quand $x$ tend vers $+ \infty$, et ceux tr�s rapidement.  Il s'agit de caract�riser un peu ce tr�s rapidement, du moins de dire que la fonction exponentielle crois bien plus vite que $x^n$. Pour cela on va �tudier la forme ind�termin�e $\limpinf \dfrac{e^x}{x^n} $. 

\begin{thm}
$$ \text{Pour tout} \; n \in \mathbb{N}, \;\;\limpinf \dfrac{e^x}{x^n}= +\infty$$
\end{thm} 

\begin{demo}
Pour tout entier naturel $n$, on pose la fonction $f_n: x \mapsto \dfrac{e^x}{x^{n+1}}$, d�finie sur $]0;+ \infty[$. \\

$f_n$ est d�rivable sur $]0;+ \infty[$ et $f_n'(x)=\dfrac{e^x x^{n+1}-(n+1)e^x x^n} {x^{2n+1}}$ et l'on a  $f'_n(x)=\dfrac{(x-n-1)e^x}{x^{n+2}}$

On remarque que: 

\begin{itemize}
\item[$\bullet$] $f_n'(x)$ s'annule pour $x=n+1$ et donc $f_n(n+1)=\dfrac{e^{n+1}}{(n+1)^{n+1}}$ est donc soit un maximum ou un minimum de la fonction $f_n$.
\item[$\bullet$] Si $x\in ]0; n+1[$ alors $f_n'(x)<0$ donc $f_n$ est strictement d�croissante sur $ [0; n+1]$
\item[$\bullet$] Si $x\in ]n+1;+ \infty$ alors $f_n'(x)>0$ donc $f_n$ est strictement croissante sur $ [n+1;+ \infty[$.\\
\end{itemize}

On d�duit alors que $f_n$ admet un minimum en $x=n+1$.\\

Ainsi pour tout $x>0$, $\dfrac{e^{n+1}}{(n+1)^{n+1}} \ie \dfrac{e^x}{x^{n+1}}$. Puisque $x>0$ on peut multiplier cette in�galit� par $x$ en conservant l'ordre, d'o�

$$\dfrac{xe^{n+1}}{(n+1)^{n+1}} \ie \dfrac{e^x}{x^{n}}$$

Or $\limpinf \dfrac{xe^{n+1}}{(n+1)^{n+1}} = + \infty$ donc on a n�c�ssairement $\limpinf  \dfrac{e^x}{x^{n}} = + \infty$ d'apr�s la derni�re in�galit�.

\end{demo}


\end{document}